Obsérvese que este razonamiento se aplica siempre que se tiene una relación
de equivalencia (o incluso solo una preequivalencia, \cf \cite{VIII.3}) en un
preesquema~$X'$, 
\ifthenelse{\boolean{orig}}
{finita y localmente libre;}
{finita y localmente libre;}
por lo demás, \Ref{VIII.7.6} es también un caso particular del resultado
análogo para las preequivalencias finitas y localmente libres, \cf \loccit
La misma observación vale para~\Ref{VIII.7.7} más abajo.

Una vez obtenida la existencia de un abierto cuasiafín saturado~$U'$ que
contiene~$x'$, también puede recurrirse a~\Ref{VIII.7.9}
y~\Ref{VIII.7.2}, con lo cual se evita el uso de las normas.

Obsérvese además que, en las condiciones de~\Ref{VIII.7.6}, si el dato de
descenso sobre~$X'$ es efectivo y~$X'$ proviene de~$X$ sobre~$S$, entonces
el morfismo~$X'\to X$ es finito, localmente libre y sobreyectivo, pues se
deduce de~$g$ por el cambio de base~$X\to S$. Se sigue (EGA~II 6.6.4) que,
si~$X'$ es cuasiproyectivo sobre~$S'$ y, por tanto, sobre~$S$, entonces~$X$
es cuasiproyectivo sobre~$S$ (se obtiene un haz invertible relativamente
amplio sobre~$X$ tomando la \emph{norma} de un haz invertible sobre~$X'$
relativamente amplio sobre~$S$, o sobre~$S'$, lo que viene a ser lo mismo).
Se obtiene así:

\begin{corollaire}
\label{VIII.7.7}
Un morfismo~$g\colon S'\to S$ finito localmente libre y sobreyectivo es un
morfismo de descenso efectivo para la categoría fibrada de los preesquemas
cuasiproyectivos sobre otros; \ie para todo~$X'$ cuasiproyectivo sobre~$S'$,
todo dato de descenso sobre~$X'$ con respecto a~$g$ es efectivo, y el
$S$-preesquema descendido~$X$ es cuasiproyectivo sobre~$S$.
\end{corollaire}

\begin{proposition}
\label{VIII.7.8}
Sea~$g\colon S'\to S$ un morfismo fielmente plano y cuasicompacto. Entonces
~$g$ es un morfismo de descenso efectivo para la categoría fibrada de los
preesquemas~$Z$ cuasicompactos sobre un preesquema~$T$, provistos de un haz
invertible amplio
\marginpar{225}
con respecto a~$T$. En particular, para todo preesquema~$X'$ sobre~$S'$,
provisto de un dato de descenso con respecto a~$g\colon S'\to S$, y todo haz
invertible~$\mathcal{L}'$ sobre~$X'$ amplio con respecto a~$S'$, provisto
también de un dato de descenso con respecto al dado sobre~$X'$ (\ie provisto
de un isomorfismo entre~$q_{1}^*(\mathcal{L}')$ y~$q_{2}^*(\mathcal{L}')$ que
satisface la condición habitual de transitividad). Entonces el dato de
descenso sobre~$X'$ es efectivo, y el haz invertible~$\mathcal{L}$ sobre el
preesquema descendido~$X$, deducido de~$\mathcal{L}'$ por descenso, es
amplio con respecto a~$S$.
\end{proposition}

La demostración es enteramente análoga a la de~\Ref{VIII.5.8}. Obsérvese
que, en el álgebra graduada cuasicoherente~$\mathcal{S}'$ sobre~$S'$
definida por~$\mathcal{L}'$, hay un dato de descenso que permite construir,
gracias a~\Ref{VIII.1.1}, un álgebra graduada
cuasicoherente~$\mathcal{S}$ sobre~$S$, de donde resulta un
~$P=\Proj(\mathcal{S})$ sobre~$S$ tal que~$P'=\Proj(\mathcal{S}')$, junto con
su dato de descenso, se identifica con~$P\times_{S}S'$. Como, por hipótesis,
~$X'$ se identifica con un abierto de~$P'$, necesariamente estable por el
dato de descenso sobre~$P'$, el dato de descenso sobre~$X'$ es también
efectivo y se obtiene el preesquema descendido como un abierto de~$P$. Se
deja el detalle al lector. --- En particular, tomando
~$\mathcal{L}'=\mathcal{O}_{X'}$, se obtiene:

\begin{corollaire}
\label{VIII.7.9}
Sea~$g\colon S'\to S$ un morfismo fielmente plano y cuasicompacto, y sea
~$X'$ un preesquema \emph{cuasiafín} sobre~$S'$. Entonces todo dato de
descenso sobre~$X'$ con respecto a~$g$ es efectivo, y el preesquema
descendido~$X$ es cuasiafín sobre~$S$.
\end{corollaire}

En virtud de~\Ref{VIII.6.2}, este resultado se aplica, en particular, si~$S'$
es localmente noetheriano y~$X'$ es cuasifinito y separado sobre~$S'$; más
generalmente, si~$S'$ es cualquiera y~$X'$ es de presentación finita,
cuasifinito y separado sobre~$S'$ (\cf \Ref{VIII.6.6}).

\begin{remarques}
\label{VIII.7.10}
Los resultados dados en este número agotan los criterios de efectividad
conocidos actualmente, y probablemente incluso los criterios útiles
existentes\footnote{Esta opinión resultó parcialmente errónea; véanse, por
ejemplo, J.P\ptbl \textsc{Murre}, Sém. Bourbaki 294 (Apéndice), mayo de 1965,
y resultados especiales (en particular, de \textsc{Néron} y
\textsc{Raynaud}). Para el descenso de los esquemas en grupos, \cf
M\ptbl Raynaud, tesis (1968).}. Obsérvense los contraejemplos siguientes en
apoyo de esta afirmación:
\begin{enumerate}
\item[(i)]
\marginpar{226}
Si~$S$ es el espectro de un cuerpo y~$S'$ es el espectro de una extensión
cuadrática galoisiana, puede encontrarse un~$X'$ sobre~$S'$, propio y liso
sobre~$S'$, de dimensión 3, provisto de un dato de descenso que no es
efectivo (Serre).
\item[(ii)] Puede encontrarse un~$S$ que sea el espectro de un anillo local regular de
dimensión 3 (si se quiere, el anillo local de un esquema algebraico sobre un
cuerpo de característica dada), un recubrimiento principal~$T$ de~$S$ con
grupo~$\ZZ/2\ZZ$, tal que, si~$t$ designa uno de los puntos de~$T$ situados
sobre el punto cerrado~$s$ de~$S$, y
\ifthenelse{\boolean{orig}}
{$S'=T-s$,}
{$S'=T-\{t\}$,}
pueda encontrarse un~$X'$ \emph{proyectivo} sobre~$S'$, regular, provisto de
un dato de descenso con respecto a~$g\colon S'\to S$, sin que este dato de
descenso sea efectivo.
\end{enumerate}

Para estas construcciones se utiliza el ejemplo de Hironaka de variedades
no proyectivas. Para~(i), basta utilizar el hecho de que puede encontrarse
sobre~$k$ un esquema propio y liso~$X_{0}$ de dimensión 3 sobre el cual
~$G=\ZZ/2\ZZ$ actúa sin inercia y en el que existen dos puntos~$a,b$
racionales
\ifthenelse{\boolean{orig}}
{sobre~$P$,}
{sobre~$S$,}
conjugados por la acción de~$G$, que no están contenidos en un abierto afín. Se pone
entonces~$X'=X_{0}\times_{k}k'$ y se hace actuar~$G$ sobre~$X'$ mediante las
acciones de~$G$ sobre los dos factores, lo que da un dato de descenso
sobre~$X'$ con respecto a~$g\colon \Spec(k')\to \Spec(k)$. Por encima de~$a$
\resp de~$b$, hay exactamente un punto~$a'$ \resp $b'$ (con extensión residual
cuadrática), y~$a'$ y~$b'$ son conjugados por la acción de~$G$, pues~$X'\to X_{0}$ es
compatible con las acciones de~$G$. Entonces~$a'$ y~$b'$ no pueden estar
contenidos en un abierto afín~$U'$, pues, de lo contrario,
~$U=X_{0}-\Im(X'-U')$ sería un abierto de~$X_{0}$ que contendría~$(a,b)$ y
cuya imagen inversa en~$X'$ estaría contenida en~$U'$ y sería, por tanto,
cuasiafín; luego~$U$ sería cuasiafín y, en consecuencia,~$(a,b)$ tendría un
entorno afín en~$U$.

Para~(ii), se utiliza el hecho de que, en el ejemplo de Hironaka, $X_{0}$ se
obtiene como preesquema propio sobre un~$k$-esquema proyectivo~$Y$, liso
sobre~$k$ (el morfismo~$f\colon X_{0}\to Y$ es, por lo demás, birracional,
pero poco importa); el grupo~$G$ actúa asimismo sobre~$Y$ de manera
compatible con sus acciones sobre~$X_{0}$; finalmente, poniendo
\ifthenelse{\boolean{orig}}
{$S'=Y-f(b)$,}
{$S'=Y-\{f(b)\}$,}
~$X'=X_{0}|S'$, $X'$ es proyectivo sobre~$S'$. Entonces~$X_{0}$ está
provisto de un dato de descenso natural con respecto al morfismo
\marginpar{227}
canónico~$Y\to S=Y/G$, gracias a las acciones de~$G$ sobre~$X_{0}$
compatibles con sus acciones sobre~$Y$. Este dato de descenso no es efectivo,
pues~$(a,b)$ no está contenido en un abierto afín.
\ifthenelse{\boolean{orig}}
{El dato de descenso inducido}
{El dato de descenso inducido}
sobre~$X'$ con respecto a~$g\colon S'\to S$ tampoco es entonces efectivo,
como se comprueba fácilmente.
\end{remarques}

